\documentclass[11pt]{article}
\usepackage[a4paper,margin=25mm]{geometry}
\usepackage[T1]{fontenc}
\usepackage{lmodern}
\usepackage{microtype}
\usepackage{amsmath,amssymb,amsthm,mathtools}
\usepackage{xurl}
\usepackage{needspace}
\usepackage[colorlinks=true,linkcolor=blue,citecolor=blue,urlcolor=blue]{hyperref}
\hypersetup{
  pdftitle={A Seginer theorem for random tensors with arbitrary iid integrable entries},
  pdfsubject={TR-21: expected injective norm and Euclidean fiber norms},
  pdfkeywords={random tensors, Seginer theorem, injective norm, discrepancy, TR-21},
  bookmarksnumbered=true
}
\urlstyle{same}
\linespread{1.045}
\setlength{\parskip}{0.3em plus 0.1em minus 0.1em}
\setlength{\parindent}{1.2em}
\setlength{\emergencystretch}{1.5em}
\setcounter{tocdepth}{1}
\newtheoremstyle{spaced}{9pt}{9pt}{\itshape}{}{\bfseries}{.}{0.5em}{}
\theoremstyle{spaced}
\newtheorem{theorem}{Theorem}[section]
\newtheorem{lemma}[theorem]{Lemma}
\newtheorem{proposition}[theorem]{Proposition}
\newtheorem{corollary}[theorem]{Corollary}
\theoremstyle{remark}
\newtheorem{remark}[theorem]{Remark}
\numberwithin{equation}{section}
\newcommand{\E}{\mathbb E}
\newcommand{\Pp}{\mathbb P}
\newcommand{\inj}{\mathrm{inj}}
\newcommand{\NN}{\mathbb N_0}
\newcommand{\ind}{\mathbf 1}
\newcommand{\step}[1]{\par\addvspace{0.5\baselineskip}\noindent\textbf{#1}\enspace}
\title{A Seginer theorem for random tensors\\with arbitrary iid integrable entries}
\author{}
\date{30 September 2026}

\begin{document}
\maketitle

\begin{abstract}
For every fixed order $k\ge3$, we prove that the expected injective norm
of a tensor with iid mean-zero integrable entries is comparable to its
largest expected Euclidean fiber norm, uniformly over all rectangular
dimensions and entry laws. This proves the conjecture of Lucca and
Pesenti recorded as TR-21. The proof extends Zhou--Zhu's shared-class
discrepancy method across magnitude levels and combines it with
Lucca--Pesenti's sparse-tail theorem and Seginer's matrix comparison.
\end{abstract}

\section{Main theorem and attribution}\label{sec:main}

For a real tensor $T\in\mathbb R^{n_1\times\cdots\times n_k}$, its
injective norm is
\begin{equation}\label{eq:injective-definition}
 \|T\|_{\inj}
 =\sup_{\substack{\|x_j\|_2=1\\1\le j\le k}}
 \left|\sum_{i_1=1}^{n_1}\cdots\sum_{i_k=1}^{n_k}
 T_{i_1,\ldots,i_k}\prod_{j=1}^k x_{j,i_j}\right|.
\end{equation}
Let $\mathcal F_j$ denote the one-dimensional fibers in mode $j$:
all indices except the $j$th are fixed. Define
\begin{equation}\label{eq:fiber-definition}
 F(T)=\max_{1\le j\le k}\E\max_{f\in\mathcal F_j}\|T_f\|_2.
\end{equation}
The maximum over modes in this definition is \emph{outside} the
expectation.

\begin{theorem}[Seginer comparison for iid tensors]\label{thm:main}
Fix $k\ge3$. There is a finite constant $C_k$, depending only on $k$,
such that every tensor $T\in\mathbb R^{n_1\times\cdots\times n_k}$,
$n_j\ge2$, with iid real mean-zero integrable entries satisfies
\begin{equation}\label{eq:main-comparison}
 F(T)\le\E\|T\|_{\inj}\le C_kF(T).
\end{equation}
The constant is independent of the dimensions and the common entry law;
in particular, the law may depend on the dimensions. No finite-variance
or higher-moment hypothesis is imposed.
\end{theorem}

Both expectations are finite: the injective norm and every fiber norm
are bounded by the sum of the absolute values of the entries.
The lower comparison holds pointwise with constant one. The upper
comparison is proved in Sections~\ref{sec:discrepancy}--\ref{sec:assembly}.
It answers the full integrable-law formulation of TR-21~\cite{TR} and
proves Lucca--Pesenti's Conjecture~1.8~\cite{LP}.

\begin{corollary}[Bernoulli entries at every density]\label{cor:bernoulli}
Let $X\in\mathbb R^{n_1\times\cdots\times n_k}$ have iid
$\operatorname{Bernoulli}(p)$ entries, where $k\ge3$, $n_j\ge2$, and
$p\in[0,1]$. For $W=X-\E X$,
\[
 F(W)\le\E\|W\|_{\inj}\le C_kF(W).
\]
\end{corollary}

This is immediate from Theorem~\ref{thm:main}. A separate Bernoulli
argument is retained in Appendix~\ref{app:bernoulli}; it also identifies
a summation refinement of Zhou--Zhu's deterministic heavy-tuple estimate.

\subsection{Relation to prior work and the additional step}

Seginer proved the corresponding comparison for iid mean-zero matrix
entries~\cite[Corollary~2.2]{Seg}. Lucca and Pesenti formulated its tensor
extension and proved a comparison under moment-equivalence hypotheses
\cite[Assumption~1.6, Theorem~1.7, and Conjecture~1.8]{LP}.
Their sparse-tail estimate, used below, is a separate result from that
moment-restricted comparison.

The light--heavy decomposition goes back to the Kahn--Szemer\'edi
method~\cite{FKS}. Zhou and Zhu developed this approach for sparse tensors
\cite{ZZ21} and subsequently introduced a simultaneous multibox
discrepancy estimate that charges each shared coordinate class only once
\cite[Lemmas~4.4--4.5, Proposition~4.6, and Appendix~A]{ZZ}.
Their Lemma~5.1 supplies the joint-rate argument for dependent occupancy
counts. Their stated Bernoulli concentration theorem assumes
$np\ge c\log n$~\cite[Theorem~2.1]{ZZ}.

The additional step here is a \emph{magnitude-labelled} discrepancy
estimate (Lemma~\ref{lem:weighted-discrepancy}) and a simultaneous
heavy-sum estimate (Proposition~\ref{prop:heavy}). Coordinate classes are
charged once across all magnitude labels, while fiber and variance
contributions are summed with squared amplitudes. These are extensions
of the cited Zhou--Zhu arguments, not independent origins of the
discrepancy method or its entropy bookkeeping. The proof then uses a
deterministic scale grid, a quantile truncation, and symmetrization to
obtain the full-law comparison. The matrix and rare-tail estimates are
imported in the precise forms stated next.

\subsection{External estimates used in the proof}\label{sec:inputs}

\step{Seginer's matrix comparison.}
For an iid mean-zero integrable random matrix $A$, there is a numerical
constant $C_S$ such that
\begin{equation}\label{eq:seginer-input}
 \E\|A\|_{\mathrm{op}}
 \le C_S\left(
 \E\max_i\|A_{i,:}\|_2+\E\max_j\|A_{:,j}\|_2\right).
\end{equation}
This is \cite[Corollary~2.2, with moment exponent one]{Seg}. It is used
only for the elongated rectangular regime in
Section~\ref{sec:rectangular}.

\step{Lucca--Pesenti's rare-tail estimate.}
Fix $k\ge3$. For independent integrable entries $U_i$ on $[n]^k$,
$n\ge2$, such that $\Pp(U_i\ne0)\le n^{-5/4}$ for every $i$,
\begin{equation}\label{eq:LP-input}
 \E\|U-\E U\|_{\inj}\le A_k\E\max_i|U_i|,
\end{equation}
where $A_k$ depends only on $k$. This is
\cite[Theorem~1.4 and Corollary~4.3]{LP}, with
$\varepsilon=1/4$. Neither identical distributions nor higher moments
are required, so deterministic zero padding is permitted.

The two estimates above are not reproved. The discrepancy and heavy-sum
adaptations are proved below; the elementary net and light-part
arguments follow the framework of \cite[Lemmas~4.1--4.2]{ZZ}.

\subsection{Notation and proof outline}

Write $[n]=\{1,\ldots,n\}$ and $\NN=\{0,1,2,\ldots\}$.
All logarithms are natural unless a base is shown, and
$\log_+t=\max\{0,\log t\}$. The letters $C_k$ and $C_{k,R}$ may
denote different constants at different occurrences, with only the
indicated dependence. For $\mu>0$ and $v\ge0$, define the upper
Poisson rate
\begin{equation}\label{eq:rate-definition}
 I_\mu(v)=
 \begin{cases}
 v\log(v/\mu)-v+\mu,&v\ge\mu,\\
 0,&v<\mu.
 \end{cases}
\end{equation}
For fixed $v$, this rate is nonincreasing in the reference mean $\mu$.

Padding a rectangular tensor with zeros into $[n]^k$, where
$n=\max_jn_j$, preserves its injective norm and all modewise fiber
maxima. Sections~\ref{sec:discrepancy}--\ref{sec:bounded} work with
independent bounded entries that are identically distributed on the
original rectangle and zero outside it. Section~\ref{sec:assembly}
combines that bounded-law estimate with \eqref{eq:LP-input}, handles
elongated formats using \eqref{eq:seginer-input}, and removes symmetry.
Section~\ref{sec:magnitudes} explains why the iid placement of magnitudes
cannot be replaced by an arbitrary deterministic coefficient tensor.

\section{Magnitude-labelled discrepancy}\label{sec:discrepancy}

Fix $k\ge3$ and $n\ge2$. Consider independent entries $Y_i$ on $[n]^k$, iid on an original
rectangular subarray and zero elsewhere, with $|Y_i|\le1$.
For $m\in\NN$ put $v_m=2^{-m}$ and
$p_m=\Pp(v_m/2<|Y_i|\le v_m)$ on the original subarray.
Labels with $p_m=0$ can be discarded almost surely. In each mode choose
pairwise disjoint, nonempty coordinate classes $D_j^\ell\subset[n]$,
labelled by $\ell\in\NN$. Write
$Q_b=\prod_jD_j^{\ell_j}$ for $b=(\ell_1,\ldots,\ell_k)$ and define
\[
 e_{b,m}=\#\{i\in Q_b:v_m/2<|Y_i|\le v_m\},\qquad
 \mu_{b,m}=p_m|Q_b|,\qquad I_{b,m}=I_{\mu_{b,m}}(e_{b,m}).
\]
The reference mean $\mu_{b,m}$ may exceed the actual mean after padding.
For a finite nonempty collection $\mathcal B$ of pairs $(b,m)$, let
\begin{align*}
 H(\mathcal B)&=\sum_{j=1}^k\sum_{\ell\text{ used in mode }j}h(D_j^\ell),
 &h(S)&=|S|\log\frac{en}{|S|},\\
 c(b,m)&=\log2+2\sum_{j=1}^k\log(\ell_j+2)+2\log(m+2).
\end{align*}
Each used coordinate class contributes once to $H$, irrespective of the
number of its coordinate-box or magnitude-label appearances. Set
$H(\varnothing)=0$; sums over an empty collection are zero.

\begin{lemma}[Joint discrepancy with magnitude labels]\label{lem:weighted-discrepancy}
For $R>0$ and $K=2(R+k+3)$, with probability at least $1-n^{-R}$,
simultaneously for every such assignment and finite collection,
\begin{equation}\label{eq:weighted-discrepancy}
 \sum_{(b,m)\in\mathcal B}I_{b,m}
 \le K\left(H(\mathcal B)+\sum_{(b,m)\in\mathcal B}c(b,m)\right).
\end{equation}
\end{lemma}
\begin{proof}
We adapt the shared-class union bound of \cite[Lemmas~4.4--4.5]{ZZ}
and the joint-rate argument of \cite[Lemma~5.1]{ZZ} to box--magnitude pairs.
Fix a valid assignment and collection. An entry belongs to at most one
selected pair $(b,m)$: distinct coordinate boxes are disjoint and an
entry has at most one magnitude level. For nonnegative $t_{b,m}$,
independence across entries and $1+u\le e^u$ therefore give
\[
 \E\exp\left(\sum_{\mathcal B}t_{b,m}e_{b,m}\right)
 \le\exp\left(\sum_{\mathcal B}\mu_{b,m}(e^{t_{b,m}}-1)\right).
\]
The counts at different magnitude levels need not be independent.
The displayed estimate uses their exclusivity within each entry instead.
For each subset $A\subseteq\mathcal B$ and $s_a\ge0$, choose
$u_a\ge\mu_a$ with $I_{\mu_a}(u_a)=s_a$ and apply the estimate with
$t_a=\log(u_a/\mu_a)$, zero outside $A$. Markov's inequality yields
\[
 \Pp\{I_a>s_a\text{ for all }a\in A\}\le e^{-\sum_{a\in A}s_a}.
\]
Expanding the product using
$e^{z/2}-1=\frac12\int_0^\infty e^{s/2}\mathbf1_{\{z>s\}}\,ds$
and integrating this joint tail gives
$\E\exp(\frac12\sum_{\mathcal B}I_a)\le2^{|\mathcal B|}$.
The probability that \eqref{eq:weighted-discrepancy} fails for this assignment
is thus at most
\[
 2^{|\mathcal B|}e^{-KH(\mathcal B)/2}
       \prod_{a\in\mathcal B}e^{-Kc(a)/2}.
\]
Put $Z_n=\sum_{\varnothing\ne S\subset[n]}e^{-Kh(S)/2}$.
The bound $\binom ns\le(en/s)^s$ and monotonicity of
$s\log(en/s)$ on $[1,n]$ give
\[
 Z_n\le\sum_{s=1}^n(en/s)^{-(K/2-1)s}
 \le n(en)^{-(K/2-1)}\le n^{-R-1}\le1.
\]
For fixed $\mathcal B$, summing $e^{-KH/2}$ over assignments is at most
$Z_n^q\le Z_n$, where $q\ge1$ is the number of used coordinate
classes; dropping their disjointness restrictions only increases the sum.
There is no coordinate-set choice for a magnitude label.
The remaining union over finite nonempty collections is at most
\[
 Z_n\left[\prod_{b,m}(1+2e^{-Kc(b,m)/2})-1\right]
 \le Z_n(e^{S_K}-1),\qquad
 S_K=2^{1-K/2}\left(\sum_{q\ge0}(q+2)^{-K}\right)^{k+1}.
\]
Since $K\ge12$, $\sum_{q\ge0}(q+2)^{-K}\le1$ and
$S_K\le2^{1-K/2}\le1/32$, so $e^{S_K}-1<1$.
The countable union, or equivalently its increasing finite truncations,
has probability at most $Z_n\le n^{-R}$.
\end{proof}

\section{All magnitude levels in one heavy sum}\label{sec:heavy}

The three-class decomposition below follows
\cite[Proposition~4.6]{ZZ}. Its modification is to sum all magnitude
levels within each fiber, variance, and shared-class entropy estimate.

\begin{proposition}[Simultaneous heavy-sum estimate]\label{prop:heavy}
Let $Y$ be as in Section~\ref{sec:discrepancy}, and suppose that
\eqref{eq:weighted-discrepancy} holds with constant $K$ for all class
assignments and finite collections. Write
\[
 \sigma^2=\E Y_i^2\quad\text{on the original rectangle},
 \qquad \Phi=\max_{j,f\in\mathcal F_j}\|Y_f\|_2.
\]
For every $D\ge\max\{1,n\sigma^2,4\Phi^2\}$ and every tuple of unit
vectors $x_1,\ldots,x_k$, put
$q_i=\prod_j|x_{j,i_j}|$ and $\tau=\sqrt D/n$. Then
\begin{equation}\label{eq:heavy}
 \sum_{i:\,|Y_i|q_i>\tau}|Y_i|q_i
 \le\left(k4^{k-1}+e^24^{k+1}+\frac{28k}{3}K\right)\sqrt D+KB_k,
\end{equation}
where $B_k=2^k[2\log2+8(k+1)]$.
\end{proposition}

\begin{proof}
If $\tau\ge1$ there are no heavy tuples. Otherwise fix unit vectors,
write $q_i=\prod_j|x_{j,i_j}|$, and let
$D_j^\ell=\{i:2^{-\ell-1}<|x_{j,i}|\le2^{-\ell}\}$.
Empty classes are omitted. With $a_j=2^{-\ell_j}$, $s_j=|D_j^{\ell_j}|$,
$\alpha_{j,\ell}=|D_j^\ell|2^{-2\ell}$, and $a_b=\prod_ja_j$, one has
\[
 \sum_\ell\alpha_{j,\ell}\le4,
 \qquad |Q_b|a_b^2=\prod_j\alpha_{j,\ell_j}.
\]
Every tuple with $|Y_i|q_i>\tau$ belongs to a pair with
$\rho_{b,m}=na_bv_m/\sqrt D>1$. Since $D\ge1$, such pairs satisfy
$m+\sum_j\ell_j<\log_2n$ and form a finite collection.
For these pairs put $\lambda_{b,m}=e_{b,m}/\mu_{b,m}$.
The nonnegative heavy sum is at most $\sum a_bv_me_{b,m}$.
Let $\mathcal B$ denote these pairs, and split them into
\begin{align*}
 \mathcal C_1&=\{(b,m)\in\mathcal B:na_j^2\le\rho_{b,m}
                   \text{ for some }j\},\\
 \mathcal C_2&=\{(b,m)\in\mathcal B:\lambda_{b,m}\le e^2\rho_{b,m}\}
                   \setminus\mathcal C_1,\\
 \mathcal C_3&=\mathcal B\setminus(\mathcal C_1\cup\mathcal C_2).
\end{align*}
\step{The squared fiber charge in $\mathcal C_1$.}
Fix $j$ and the other coordinate labels, and write $P=\prod_{i\ne j}a_i$.
The condition is $a_j\le v_mP/\sqrt D$. The sets $D_j^{\ell_j}$ are
disjoint. If $c_m(f)$ counts level-$m$ entries in fiber $f$, then
\[
 \sum_m v_m^2c_m(f)\le4\sum_{i\in f}Y_i^2\le D.
\]
It follows, summing the qualifying pairs, that
\[
 \begin{aligned}
 \sum_{\substack{\ell_j,m:\,(b,m)\in\mathcal B\\
                 a_j\le v_mP/\sqrt D}}a_jPv_me_{b,m}
 &\le\frac{P^2}{\sqrt D}
       \sum_{f\text{ in the fixed other classes}}\sum_m v_m^2c_m(f)\\
 &\le\sqrt D\prod_{i\ne j}\alpha_{i,\ell_i}.
 \end{aligned}
\]
The sum on the right controls all the qualifying pairs because the
classes $D_j^{\ell_j}$ are disjoint. Summing over the other labels
and then $j$ (overcounting a pair that qualifies for several modes is
harmless) gives
\begin{equation}\label{eq:C1}
 \sum_{\mathcal C_1}a_bv_me_{b,m}\le k4^{k-1}\sqrt D.
\end{equation}
\step{The combined variance charge in $\mathcal C_2$.}
Here
\[
 \frac{a_bv_me_{b,m}}{\sqrt D}
 \le e^2\frac{np_mv_m^2}{D}\prod_j\alpha_{j,\ell_j}.
\]
Since $\sum_mp_mv_m^2\le4\sigma^2$, summing all labels gives
\begin{equation}\label{eq:C2}
 \sum_{\mathcal C_2}a_bv_me_{b,m}\le e^24^{k+1}\sqrt D.
\end{equation}
\step{Shared entropy in $\mathcal C_3$.}
The defining inequalities are $1<\rho_{b,m}<na_j^2$ for every $j$ and
$\lambda_{b,m}>e^2\rho_{b,m}$. Thus
$I_{b,m}\ge e_{b,m}\log(e\rho_{b,m})$.
For $u\ge0$, let
\[
 \mathcal C_3(u)=\left\{(b,m)\in\mathcal C_3:
  u<\frac{a_bv_m}{\sqrt D\log(e\rho_{b,m})}\right\}.
\]
Integrating \eqref{eq:weighted-discrepancy} over these subcollections gives
\begin{equation}\label{eq:C3-integral}
 \frac1{\sqrt D}\sum_{\mathcal C_3}a_bv_me_{b,m}
 \le K\int_0^\infty H(\mathcal C_3(u))\,du
    +K\sum_{\mathcal C_3}
          \frac{c(b,m)a_bv_m}{\sqrt D\log(e\rho_{b,m})}.
\end{equation}
The interval endpoint equals $\rho_{b,m}/[n\log(e\rho_{b,m})]$.
Monotonicity of $t/\log(et)$ for $t\ge1$ bounds the lifetime of each
used class $D_j^\ell$, across \emph{all} magnitude labels, by
$a_\ell^2/\log(ena_\ell^2)$, with $na_\ell^2>1$.
Since $h(D_j^\ell)=|D_j^\ell|\log(en/|D_j^\ell|)$, its integrated cost is at most
\[
 \alpha_{j,\ell}
 \left(1+\frac{\log_+(1/\alpha_{j,\ell})}{\log(en4^{-\ell})}\right).
\]
For $t_\ell=n4^{-\ell}>1$, this is bounded by
$2\alpha_{j,\ell}+1/t_\ell$: use
$\log_+(1/\alpha)\le\log t+1/(\alpha t)$.
Since $\sum_{t_\ell>1}1/t_\ell<4/3$, the entropy integral is at most
$28k/3$. This is the entropy estimate from \cite[Appendix~A]{ZZ},
with one lifetime for each shared class across all magnitude labels.
It imposes no lower bound on $n\sigma^2$.
For the label term, $\log(e\rho)>1$ and $\log(q+2)\le q+1$ give
\[
 \sum_{b,m\ge0}c(b,m)2^{-(\sum_j\ell_j+m)}
 \le B_k:=2^k[2\log2+8(k+1)].
\]
Combining \eqref{eq:C1}--\eqref{eq:C3-integral} proves
\eqref{eq:heavy}. In particular, no number of magnitude levels enters
the bound.
\end{proof}

\section{Bounded symmetric entry laws}\label{sec:bounded}

\begin{theorem}[Bounded-law high-probability estimate]\label{thm:bounded}
Fix $k\ge3$ and $R>0$. Let the entries of $Y$ be iid and symmetric on
a rectangular subarray of $[n]^k$, $n\ge2$, and zero elsewhere.
Assume $|Y_i|\le1$, and set
\[
 \sigma^2=\E Y_i^2\quad\text{on the original rectangle},
 \qquad \Phi=\max_{j,f\in\mathcal F_j}\|Y_f\|_2.
\]
There is a constant $C_{k,R}$ such that
\begin{equation}\label{eq:bounded-hp}
 \Pp\!\left\{\|Y\|_{\inj}>
 C_{k,R}\bigl(1+\sqrt{n\sigma^2}+\Phi\bigr)\right\}\le2n^{-R}.
\end{equation}
\end{theorem}

\begin{proof}
For fixed unit vectors $x_1,\ldots,x_k$, put
$q_i=\prod_j|x_{j,i_j}|$. For deterministic
$D_0\ge\max\{1,n\sigma^2\}$, the light sum has summands
\[
 Y_i\prod_jx_{j,i_j}\,
 \mathbf1_{\{|Y_i|q_i\le\sqrt{D_0}/n\}}.
\]
They are centered because they are odd functions of symmetric entries,
independent, bounded by $\sqrt{D_0}/n$, and have total variance at most
$\sigma^2\le D_0/n$. Bernstein's inequality gives
\[
 \Pp\{|L_Y|>B\sqrt{D_0}\}
 \le2\exp\left(-\frac{nB^2}{2+2B/3}\right).
\]
Choose a deterministic $1/(2k)$-net $\mathcal T$ of the unit
sphere with $|\mathcal T|\le(1+4k)^n$, as in
\cite[Lemma~4.1]{ZZ}. Replacing one vector at a time gives
\[
 \|Y\|_{\inj}\le2\max_{(x_1,\ldots,x_k)\in\mathcal T^k}
 |Y(x_1,\ldots,x_k)|.
\]
This inequality is applied to the full multilinear form, not to its
truncated light part. Union over the product net and the grid
\[
 \mathcal D_n=\{2^q:0\le q\le\lceil\log_2(4n)\rceil\}.
\]
This grid has $O(\log(en))$ elements and maximum in $[4n,8n)$.
With $B=B_{k,R}$ sufficiently large, all allowed light bounds hold
with failure probability at most $n^{-R}$. After sampling, select the
smallest grid value
$D_0\ge\max\{1,n\sigma^2,4\Phi^2\}$.
This exists because $\sigma^2\le1$ and $\Phi^2\le n$, and it is less
than twice that maximum. Combine Lemma~\ref{lem:weighted-discrepancy},
Proposition~\ref{prop:heavy}, and the net bound.
If $D_0\ge n^2$ the heavy set is empty and the same light bound suffices.
The two exceptional events have total probability at most
$2n^{-R}$. Since $D_0$ is less than twice
$\max\{1,n\sigma^2,4\Phi^2\}$ and $D_0\ge1$, the additive $KB_k$
in \eqref{eq:heavy} is absorbed into a constant times $\sqrt{D_0}$.
This proves \eqref{eq:bounded-hp}.
\end{proof}

\begin{remark}
The grid has the fixed lower endpoint $1$. Its corresponding additive
amplitude term in \eqref{eq:bounded-hp} is retained, not discarded;
Section~\ref{sec:assembly} charges it to $F$ after truncation and rescaling.
The simultaneous deterministic grid is what permits a random final
choice of $D_0$ without conditioning away centering or independence.
\end{remark}

\clearpage
\section{Proof of the full-law comparison}\label{sec:assembly}

We first prove the upper bound for symmetric iid integrable entries,
then reduce the centered case to it by symmetrization.

\subsection{Elongated rectangular formats}\label{sec:rectangular}

Let $Z$ have symmetric iid integrable entries in the original format. Write $n=\max_jn_j$, $V=\prod_jn_j$.
If $n\ge\prod_{i\ne j}n_i$ for a longest mode $j$, flatten into an
$n\times P$ iid matrix with $P\le n$. Its maximal column norm is the
mode-$j$ fiber maximum. To compare rows and columns, couple iid squared
entries along a line of length $nP$ and partition it into blocks of
length $P$ and, separately, length $n$. Each shorter block intersects
at most two longer ones. Hence the expected maximal row norm is at
most $\sqrt2$ times the expected maximal column norm. Seginer's estimate \eqref{eq:seginer-input} and
$\|Z\|_{\inj}\le\|\mathrm{flatten}(Z)\|_{\mathrm{op}}$ give
$\E\|Z\|_{\inj}\le C_S(1+\sqrt2)F(Z)$.
The coupling compares the two marginal expectations and need not
be the row and column partition of one matrix realization.

\subsection{Quantile truncation and the rare component}\label{sec:truncation}

It remains to consider $V>n^2$. Pad the tensor with zeros to $[n]^k$,
without changing its norm or fiber maxima, and put
\[
 \theta=\inf\{t\ge0:\Pp(|Z_i|>t)\le n^{-5/4}\},\quad
 Z^{\rm rare}_i=Z_i\mathbf1_{\{|Z_i|>\theta\}},\quad
 Z^{\rm c}_i=Z_i\mathbf1_{\{|Z_i|\le\theta\}}.
\]
The quantile is finite and, by right-continuity of the strict tail,
$\Pp(|Z_i|>\theta)\le n^{-5/4}$.
Both components are symmetric, independent across entries, and
integrable. Equation \eqref{eq:LP-input} gives
\[
 \E\|Z^{\rm rare}\|_{\inj}
 \le A_k\E\max_i|Z^{\rm rare}_i|\le A_kF(Z).
\]
If $\theta=0$ this already handles $Z$. For $\theta>0$, each $0\le t<\theta$
has strict-tail probability greater than $n^{-5/4}$. Independence gives
\[
 \E\max_i|Z_i|\ge
 t\bigl(1-e^{-Vn^{-5/4}}\bigr).
\]
Letting $t\uparrow\theta$ and using $V\ge n^2$, $n\ge2$, yields
\begin{equation}\label{eq:theta-charge}
 \theta\le2\E\max_i|Z_i|\le2F(Z).
\end{equation}
\subsection{The bounded component and its variance term}

Let $\sigma_{\rm c}^2=\E(Z_i^{\rm c})^2$ and
$\Phi_{\rm c}=\max_f\|Z_f^{\rm c}\|_2$.
Apply Theorem~\ref{thm:bounded} to $Y=Z^{\rm c}/\theta$.
On its failure event use
$\|Z^{\rm c}\|_{\inj}\le\theta n^{k/2}$.
Taking $R=k+2$ gives
\begin{equation}\label{eq:bounded-expectation}
 \E\|Z^{\rm c}\|_{\inj}
 \le C_k\bigl(\theta+\sqrt{n\sigma_{\rm c}^2}+\E\Phi_{\rm c}\bigr)
       +2\theta n^{k/2-R}.
\end{equation}
Monotonicity of each squared fiber energy implies
$\E\Phi_{\rm c}\le\sum_j\E\max_{f\in\mathcal F_j}\|Z_f\|_2
\le kF(Z)$.
For the variance term put
$S=\sum_{i\text{ original}}(Z_i^{\rm c})^2$ and
$\mu=\E S=V\sigma_{\rm c}^2$.
Each summand lies in $[0,\theta^2]$, so independence gives
$\E S^2\le\mu^2+\theta^2\mu$.
If $\mu\ge8\theta^2$, the Paley--Zygmund inequality yields
\[
 \Pp(S\ge\mu/2)\ge\frac{\mu^2}{4\E S^2}
 \ge\frac1{4(1+1/8)}=\frac29.
\]
The longest mode has $V/n$ fibers; its maximal squared fiber energy
is at least $nS/V$. Therefore
\[
 F(Z)\ge\sqrt{n/V}\,\E\sqrt S
 \ge\frac{2}{9\sqrt2}\sqrt{n\sigma_{\rm c}^2}.
\]
If $\mu<8\theta^2$, then
$\sqrt{n\sigma_{\rm c}^2}<\sqrt8\,\theta\sqrt{n/V}
\le\sqrt8\,\theta$, absorbed by \eqref{eq:theta-charge}.
Every term in \eqref{eq:bounded-expectation}, including its failure contribution,
is thus at most a constant depending only on $k$ times $F(Z)$.
The triangle inequality and the rare-component estimate therefore
prove $\E\|Z\|_{\inj}\le C_kF(Z)$ for symmetric iid integrable laws.

\subsection{Symmetrization and completion of the proof}

For the centered iid tensor $T$ in Theorem \ref{thm:main}, let $T'$
be an independent copy and set $Z=T-T'$. Its entries are symmetric iid
and integrable. Conditional Jensen and $\E T'=0$ give
$\E\|T\|_{\inj}\le\E\|Z\|_{\inj}$.
For each mode, the triangle inequality applied before taking the
maximum and expectation gives $F(Z)\le2F(T)$.
The symmetric comparison therefore proves the asserted upper bound.
The lower bound is pointwise: fix all but one argument to coordinate
vectors and optimize over the remaining fiber.


Equivalently, the final reduction reads
\[
 \E\|T\|_{\inj}\le\E\|T-T'\|_{\inj}
 \le C_kF(T-T')\le2C_kF(T).
\]
Renaming the constant proves Theorem~\ref{thm:main}.

\section{Why iid magnitude placement matters}\label{sec:magnitudes}

A layerwise triangle inequality would replace the fiber norm of the
whole tensor by a sum of layerwise fiber scales. Such a sum need not be
controlled without a loss depending on the number of layers.
Zhou--Zhu's deterministic-weight bound
\cite[Corollary~2.2]{ZZ}, expressed through a maximum weight and a maximum
sampling probability, does not supply the joint estimate needed here.
The proof above instead sums squared amplitudes in
\eqref{eq:C1} and \eqref{eq:C2}, and charges each coordinate class once
across all magnitude levels in \eqref{eq:C3-integral}.

The iid placement of magnitudes cannot simply be discarded. To see
this, let $M=2^{2d^2}$ and construct an order-three coefficient tensor
$A$ supported on $M$ disjoint $d\times d$ all-one blocks, with disjoint
first- and second-mode coordinate sets and a separate third-mode
coordinate for each block. Pad with zeros into cubic format. Then
\[
 \max_f\|A_f\|_2=\sqrt d.
\]
For independent Rademacher signs $\varepsilon_i$, an individual block
of $\varepsilon\odot A$ is all positive with probability $2^{-d^2}$,
and these events are independent across blocks. Thus
\[
 \Pp\{\text{some block is all positive}\}
 =1-(1-2^{-d^2})^M\longrightarrow1.
\]
On that event, choose normalized all-one vectors on the corresponding
first- and second-mode blocks and the appropriate third-mode coordinate
vector. Their evaluation is $d$, so
\[
 \E_\varepsilon\|\varepsilon\odot A\|_{\inj}
 \ge d\bigl[1-(1-2^{-d^2})^M\bigr].
\]
Consequently, a uniform comparison with the largest fiber norm is false
for arbitrary deterministic magnitudes, even after averaging signs.
This does not contradict Theorem~\ref{thm:main}: its entries are iid,
whereas this coefficient construction fixes a highly nonuniform support.

\clearpage
\appendix
\section{A direct Bernoulli argument}\label{app:bernoulli}

This appendix proves Corollary~\ref{cor:bernoulli} directly from the
Zhou--Zhu discrepancy and heavy-tuple framework, without using
\eqref{eq:seginer-input} or \eqref{eq:LP-input}. Its additional steps
are the retained-denominator estimate below and a simultaneous choice
of the degree parameter. The all-density comparison is a deduction
here, not a theorem stated in \cite{ZZ}.

\subsection{The retained-denominator summation}

For $t>1$, define
\[
 S_k(t)=\sum_{\substack{\ell\in\NN^k\\t2^{-|\ell|_1}>1}}
 \frac{2^{-|\ell|_1}}{1+\log(t2^{-|\ell|_1})}.
\]
\begin{lemma}[A dyadic summation bound]\label{lem:summation}
For fixed $k$, $S_k(t)\leq C_k/(1+\log t)$ for every $t>1$.
\end{lemma}
\begin{proof}
Write $u=\log_2 t$. Grouping by $s=|\ell|_1$ gives
\[
 S_k(t)=\sum_{0\leq s<u}
 \binom{s+k-1}{k-1}\frac{2^{-s}}{1+(u-s)\log2}.
\]
For $s\leq u/2$, the denominator is at least $1+(\log t)/2$, and
$\sum_{s\geq0}\binom{s+k-1}{k-1}2^{-s}=2^k$.
For $s>u/2$, use the denominator bound $1$ and the geometric tail estimate
\[
 \sum_{s>u/2}\binom{s+k-1}{k-1}2^{-s}
 \leq C_k(1+u)^{k-1}2^{-u/2}
 \leq \frac{C'_k}{1+u}.
\]
The last inequality follows from the boundedness of a fixed polynomial times an exponentially decreasing function, including bounded $u$. These two bounds prove the claim.
\end{proof}

Set $N=\max_j n_j$, $L_*=1+\lfloor\log_2 N\rfloor$, and let $1\leq D\leq N$ be a parameter. For a nonnegative tensor padded with zeros into $[N]^k$, assume every fiber has sum at most $D$. Suppose that the multibox discrepancy inequality of Zhou--Zhu \cite[Lemma~4.5]{ZZ} holds with constant $K$, labels $0,\ldots,L_*-1$, and reference means $(D/N)|Q_b|$:
\begin{equation}\label{eq:bernoulli-discrepancy}
 \sum_{b\in\mathcal B}I_{(D/N)|Q_b|}(e(Q_b))
 \leq K\bigl(H(\mathcal B)+|\mathcal B|\log(eL_*)\bigr).
\end{equation}
Here $e(Q)$ is the sum in a box, $I_\mu$ is
\eqref{eq:rate-definition}, and for coordinate classes $D_j^\ell$ and
boxes $Q_b=\prod_jD_j^{\ell_j}$,
\[
 H(\mathcal B)=\sum_{j=1}^k\sum_{\ell\text{ used in mode }j}
 |D_j^\ell|\log\frac{eN}{|D_j^\ell|}.
\]
For unit vectors, use the heavy threshold $\tau=\sqrt D/N$.
The dyadic classes in \cite[Proposition~4.6]{ZZ} have amplitudes
$a_\ell=2^{-\ell}$; classes that occur in a heavy tuple require at most
$L_*$ labels. In its difficult class $\mathcal C_3$, write
$a_b=2^{-|b|_1}$ and $\rho_b=Na_b/\sqrt D>1$.
The estimate preceding the simplification in
\cite[equations~(4.15)--(4.17)]{ZZ} gives an \emph{absolute},
unnormalized label-cost contribution of at most
\[
 K\log(eL_*)\sum_{b\in\mathcal C_3}\frac{a_b}{\log(e\rho_b)}.
\]
Lemma \ref{lem:summation}, with $t=N/\sqrt D$, bounds this by
\[
 \frac{C_k K\log(eL_*)}{1+\log(N/\sqrt D)}\leq C'_k K,
\]
because $D\leq N$ implies $N/\sqrt D\geq\sqrt N$. The other terms in \cite[proof of Proposition~4.6]{ZZ} are unchanged:
equations~(4.12) and (4.13) give respectively $k4^{k-1}\sqrt D$ and $e^24^k\sqrt D$, while (4.16) gives $(28/3)kK\sqrt D$. The entropy estimate in \cite[Appendix~A]{ZZ} does not require
a lower bound proportional to $\log N$. Since $D\geq1$, the resulting heavy sum is at most $C_{k,K}\sqrt D$.

Thus the proof yields the deterministic heavy bound for $1\leq D\leq N$; the assumption $D\geq c\log N$ is unnecessary after this refinement. This does \emph{not} mean that the final bound can use $D=Np$ below logarithmic degree: the fiber-degree hypothesis must still hold.

\subsection{Probabilistic inputs with a random degree scale}
It suffices to treat $0<p\leq1/2$; complementing Bernoulli entries handles $p>1/2$, and $p=0,1$ are trivial. Pad $X$ and $\E X$ with zeros into $[N]^k$, which preserves the injective norm and fiber maxima. Put
\[
 d=Np,\qquad
 \Delta=\max_f\sum_{i\in f}X_i,\qquad
 D=\left\lceil\max\{1,d,\Delta\}\right\rceil\in\{1,\ldots,N\}.
\]
Fix a failure exponent $R>0$. The padded counts in disjoint boxes are independent binomials with means at most $p|Q|$. Their upper tails are bounded by the Poisson rate $I_{p|Q|}$: this follows either from their exponential moment bound or from monotonicity in the reference mean. Consequently, \cite[proofs of Lemmas~4.4--4.5]{ZZ} give, with probability at least $1-N^{-R}$,
\[
 \sum_{b\in\mathcal B}I_{p|Q_b|}(e(Q_b))
 \leq K\bigl(H(\mathcal B)+|\mathcal B|\log(eL_*)\bigr)
\]
for every partial partition and nonempty box collection. The zero padding therefore does not require identical distributions in this input.

For fixed $v$, $I_\mu(v)$ is nonincreasing in $\mu>0$: on $\mu<v$, its derivative is $1-v/\mu\leq0$, and it is zero when $\mu\geq v$. Since $D\geq Np$, the same event implies \eqref{eq:bernoulli-discrepancy} for every allowed $D$. Every fiber has degree at most $D$ by definition.

For the light part, take the deterministic product net of \cite[Lemma~4.1]{ZZ}. For each deterministic integer $a\in[1,N]$ with $a\geq d$, use threshold $\tau_a=\sqrt a/N$. At a fixed net point, the centered summands are bounded by $\tau_a$ and have total variance at most $p\leq a/N$. Bernstein's inequality gives
\[
 \Pp\{|L_W|>B\sqrt a\}
 \leq 2\exp\left\{-\frac{N B^2}{2+(2/3)B}\right\}.
\]
Choose $B=B_{k,R}$ sufficiently large to union bound over at most $(1+4k)^{kN}$ product-net points and at most $N$ integer parameters. The light estimate then holds simultaneously for all these parameters, outside an event of probability at most $N^{-R}$. Selecting $a=D$ is valid even though $D$ is random.

The heavy part of the mean is bounded by
\[
 p\sum_{i:\,q_i>\sqrt D/N}q_i
 \leq\frac{pN}{\sqrt D}\sum_i q_i^2
 =\frac d{\sqrt D}\leq\sqrt D,
 \qquad q_i=\prod_{j=1}^k|x_{j,i_j}|.
\]
The refined deterministic bound controls the nonnegative heavy part of $X$. Combining these estimates and the factor-two net approximation gives
\begin{equation}\label{eq:bernoulli-hp}
 \|W\|_{\inj}\leq C_{k,R}\sqrt D
 \quad\text{outside an event of probability at most }2N^{-R}.
\end{equation}
Since $|W_i|\leq1$, the norm is always at most $N^{k/2}$. Hence
\begin{equation}\label{eq:bernoulli-expectation}
 \E\|W\|_{\inj}
 \leq C_{k,R}\E\sqrt D+2N^{k/2-R}.
\end{equation}

\subsection{Charging the degree scale to the actual fibers}
Let $V=\prod_jn_j$ and $\lambda=Vp$. Write $\Delta_j$ for the maximum count in a mode-$j$ fiber. A fiber containing $m$ ones obeys
\[
 \|W_f\|_2^2=m(1-p)^2+(n_j-m)p^2\geq m/4.
\]
It follows that
\[
 \E\sqrt\Delta\leq\sum_{j=1}^k\E\sqrt{\Delta_j}\leq2kF(W).
\]
On the event of at least one nonzero entry of $X$, each mode has a centered fiber norm at least $1-p\geq1/2$. Thus
\begin{equation}\label{eq:bernoulli-presence}
 F(W)\geq\tfrac12\bigl(1-(1-p)^V\bigr)
 \geq\tfrac12(1-e^{-\lambda}).
\end{equation}
If $d=Np\geq1$, a fixed fiber in a longest mode has count $M\sim\operatorname{Bin}(N,p)$. The Paley--Zygmund inequality gives
\[
 \Pp\{M\geq d/2\}\geq\frac{d^2}{4\E M^2}
 \geq\frac{d}{4(d+1)}\geq\tfrac18.
\]
Therefore $F(W)\geq\sqrt d/(16\sqrt2)$.

If $\lambda\geq1$, \eqref{eq:bernoulli-presence} bounds $F(W)$ below by a numerical constant, and $\sqrt d\leq C F(W)$ whether $d\geq1$ or $d<1$. Since
\[
 \sqrt D\leq\sqrt2\bigl(1+\sqrt d+\sqrt\Delta\bigr),
\]
we obtain $\E\sqrt D\leq C_kF(W)$. Choose $R=k+2$ in \eqref{eq:bernoulli-expectation}; its residual is bounded by a constant and is absorbed by $F(W)$. This proves the upper bound in this regime.

If $\lambda<1$, let $M_{\rm tot}$ be the total number of ones. The elementary estimate
\[
 \E\|W\|_{\inj}
 \leq\E\|X\|_F+\|\E X\|_{\inj}
 =\E\sqrt{M_{\rm tot}}+p\sqrt V
 \leq\lambda+\lambda=2\lambda
\]
uses $\sqrt m\leq m$ for nonnegative integers. By \eqref{eq:bernoulli-presence} and
$1-e^{-\lambda}\geq(1-e^{-1})\lambda$ on $[0,1]$, $F(W)\geq c\lambda$. The upper bound follows here as well. The lower bound in Corollary~\ref{cor:bernoulli} is immediate by fixing all but one argument to coordinate vectors. This completes the proof.

\subsection{The critical density \texorpdfstring{$p=1/n$}{p=1/n}}
\label{app:critical}

In cubic format $n^k$, a fixed mode contains $n^{k-1}$ independent $\operatorname{Bin}(n,1/n)$ fiber counts. If $\Delta_j$ is their maximum, then, for fixed $k$ and $n\to\infty$,
\[
 \E\sqrt{\Delta_j}\asymp_k\sqrt{\frac{\log n}{\log\log n}}.
\]
For completeness, Chernoff's bound
\[
 \Pp\{\operatorname{Bin}(n,1/n)\geq t\}\leq(e/t)^t
\]
and a union bound, followed by summing the tail above a sufficiently
large multiple of $\log n/\log\log n$, give the upper scale. For the lower scale, for integers $1\leq t\leq n$,
\[
 \Pp\{\operatorname{Bin}(n,1/n)=t\}\geq\tfrac14t^{-t}.
\]
Choosing $t=\lfloor((k-1)/2)\log n/\log\log n\rfloor$ makes the expected number of fibers with this count tend to infinity; independence gives occurrence with probability tending to one. A centered fiber with count $m$ has squared norm
$m(1-2/n)+1/n$, which lies between $m/4$ and $m+1/n$ for $n\ge2$.
Thus the expected centered fiber maximum has the same square-root scale. Consequently, as $n\to\infty$,
\[
 \E\|X-1/n\|_{\inj}\asymp_k F(X-1/n)
 \asymp_k\sqrt{\frac{\log n}{\log\log n}}.
\]
This benchmark also follows from \cite[proof of Proposition~4.6]{ZZ}
without the retained-denominator refinement: that proof gives a heavy bound $C_{k,K}(\sqrt D+\log(eL_*))$, and $\log(eL_*)=O(\log\log n)$ is smaller than
$\E\sqrt D\asymp_k\sqrt{\log n/\log\log n}$ in this regime.
This observation uses the random-degree and expectation reductions above;
it is not a direct application of \cite[Theorem~2.1]{ZZ}.

\phantomsection
\addcontentsline{toc}{section}{References}
\begin{thebibliography}{99}

\bibitem{FKS}
J.~Friedman, J.~Kahn, and E.~Szemer\'edi.
\newblock On the second eigenvalue of random regular graphs.
\newblock In \emph{Proceedings of the Twenty-First Annual ACM Symposium
on Theory of Computing (STOC '89)}, pages 587--598, 1989.
\newblock \href{https://doi.org/10.1145/73007.73063}{doi:10.1145/73007.73063}.

\bibitem{LP}
K.~Lucca and L.~Pesenti.
\newblock \emph{Norm Bounds for Sparse Random Tensors and Spectral Gap
of Random Hypergraphs}.
\newblock arXiv:2607.07308v1, 8 July 2026.
\newblock The rare-tail input is Theorem~1.4 and Corollary~4.3;
Theorem~1.7 is the moment-restricted comparison, and Conjecture~1.8
is the full tensor comparison.
\newblock \url{https://arxiv.org/abs/2607.07308v1}.

\bibitem{TR}
\emph{TR-21: A Seginer theorem for arbitrary independent identically
distributed tensor entries}.
\newblock Open Problems in NLA, problem catalogue. Accessed 30 September 2026.
\newblock \url{https://github.com/ajt60gaibb/OpenProblemsInNLA/blob/main/tensor-computations/TR-21/README.md}.

\bibitem{Seg}
Y.~Seginer.
\newblock The expected norm of random matrices.
\newblock \emph{Combinatorics, Probability and Computing},
9(2):149--166, 2000.
\newblock Corollary~2.2.
\newblock \href{https://doi.org/10.1017/S096354830000420X}{doi:10.1017/S096354830000420X}.

\bibitem{ZZ21}
Z.~Zhou and Y.~Zhu.
\newblock Sparse random tensors: Concentration, regularization and applications.
\newblock \emph{Electronic Journal of Statistics}, 15(1):2483--2516, 2021.
\newblock \href{https://doi.org/10.1214/21-EJS1838}{doi:10.1214/21-EJS1838}.

\bibitem{ZZ}
Z.~Zhou and Y.~Zhu.
\newblock \emph{Sharp spectral norm concentration of sparse random tensors}.
\newblock arXiv:2609.20520v1, 17 September 2026.
\newblock The proof adaptations use Lemmas~4.1--4.5 and~5.1,
Proposition~4.6, and Appendix~A. All corresponding external equation
numbers in this manuscript refer to this version.
\newblock \url{https://arxiv.org/abs/2609.20520v1}.

\end{thebibliography}
\end{document}
